\chapter{INTRODUCTION}\label{ch:intro}

\begin{doublespacing}

In the minutes and hours after a building collapses, collapsed-structure
entry is dangerous for human rescuers and a high-stakes information
problem: a survivor is somewhere in a rubble field and the probability
that they are in any particular void pocket is unevenly distributed.
Small autonomous robots---aerial drones, crawling ground units, or
tethered probes---can be deployed into the structure to search. When
several robots are available, the question becomes not merely
\emph{where should a robot look} but \emph{how should a team of robots
divide the search}. Poor coordination causes robots to check the same
locations and to burn their finite batteries before checking the most
likely ones. Good coordination \emph{on the right mathematical
foundations} can cut search distance dramatically, and it is those
foundations that this thesis studies.

This thesis addresses the following question. Given a team of $R$
battery-powered robots, a set of $n$ candidate survivor locations with
non-uniform prior probabilities, and a worst-case metric of efficiency,
what combination of communication assumptions, bid functions, and
routing decisions minimizes the distance a robot travels before finding
the target? The answer, which this thesis develops formally and
validates numerically, is that the bottleneck is not fleet size or
battery capacity in isolation but the bid function that robots use to
decide which site to visit next, together with their ability to chain
several sites into a single energy-feasible tour.

\section{Motivation}

Urban search-and-rescue (USAR) robotics has matured from tele-operated
single-unit deployments toward teams of semi-autonomous robots. Field
reports from deployments at the World Trade Center~\cite{Casper2003},
the Fukushima Daiichi plant, and several recent earthquakes consistently
identify two practical pain points: robots exhaust their batteries
before finishing their search, and multiple robots redundantly cover
the same ground when communication links drop. The two problems are
related. Under strict energy budgets, each wasted sortie---that is,
each excursion that checks a site of low prior probability or revisits
a site that a teammate already cleared---is a larger fraction of the
total searchable volume. Coordination quality therefore matters more
as energy grows scarcer, not less, which inverts the intuition that
a ``self-sufficient'' robot with enough battery can simply brute-force
the problem.

From a theoretical standpoint, multi-robot search sits at the
intersection of three literatures that have mostly evolved in
isolation. Competitive-ratio analysis, originating from the
cow-path problem~\cite{BaezaYates1993}, gives precise guarantees on
online search on simple topologies but assumes a single searcher on a
half-line. Market-based auction coordination~\cite{Dias2006} gives
practical scalability but traditionally assumes unlimited mobility and
therefore ignores energy. Energy-constrained exploration~\cite{Betke1995}
studies fuel-limited traversal but typically without auction
coordination. This thesis unifies all three into a single framework:
competitive-ratio upper bounds on auction-coordinated search under
finite energy.

\section{Scope and Contributions}

The scope of this work is algorithmic and simulation-based. Four
algorithmic models that span the energy--communication design space
are formulated, closed-form upper bounds on their expected finder
competitive ratio are derived, and the bounds are validated by 500
Monte Carlo trials per configuration on identical instances with
sweeps over the energy budget $E$ and the fleet size $R$. Physical
robot deployment is not in scope and is left as future work.

The specific contributions of this thesis are as follows.

\begin{enumerate}
\item Formalization of four search models (M1 random, M2 unconstrained
auction, M3 single-sortie auction, M4/M4* multi-sortie auction) that
systematically vary the two operational dimensions---energy and
communication---and derivation of an expected-FCR upper bound for each.
\item An ablation on the Hungarian centralized assignment baseline that
disentangles two previously conflated effects: probability-aware bidding
versus centralized allocation. Probability-awareness is shown to be the
dominant factor.
\item Identification of a bid function $p_i/d^2$ that empirically
outperforms the standard $p_i/d$ bid under energy constraints, together
with a cost-foreclosure argument that justifies the quadratic distance
penalty.
\item Validation of all four bounds across 500 trials, with parameter
sweeps over $E \in [8, 30]$ and $R \in [1, 6]$ showing the stability of
the ranking and the diminishing-returns behavior of fleet size.
\end{enumerate}

\section{Thesis Organization}

The remainder of the thesis is organized as follows. Chapter~II reviews
the related work in competitive search theory, auction-based
coordination, and probabilistic and energy-constrained search, and
positions the present contribution relative to each. Chapter~III states
the formal problem, defines the finder competitive ratio, and
introduces the sortie model and the geometric constants used in the
bounds. Chapter~IV presents the four algorithmic models and the
Hungarian baselines, together with the theoretical upper bounds on the
expected FCR. Chapter~V describes the experimental setup and reports
the main comparison results, the bounds verification, and the
parameter sweeps, and discusses the filtering thresholds and the
simplifying assumptions. Chapter~VI concludes with a summary of
findings and a discussion of future work.

\end{doublespacing}
